\centerline{\textbf{\large Scalable and Decentralized Urban Traffic Optimization}}
\centerline{\textbf{\large with Verifiable Storage for Smart City Deployments}}

\vspace{2cm}

\rightline{Autores: Kevin Mathias Galeano Saldivar, María José Duarte Kowalewski}
\vspace{0.5cm}
\rightline{Asesor: Cristian Cappo}

\vspace{0.5cm}

\noindent
\textbf{RESUMEN}

\vspace{0.5cm}

\noindent
Este trabajo presenta el diseño, la implementación y la evaluación experimental de un sistema
modular y descentralizado para la optimización adaptativa del control semafórico urbano y la
persistencia verificable de los datos asociados. La arquitectura propuesta integra cuatro
microservicios: un módulo de simulación de tráfico basado en SUMO (traffic-sim),
un módulo de análisis y optimización que combina inferencia difusa tipo Mamdani con
optimización por enjambre de partículas (traffic-sync), un módulo de almacenamiento
distribuido basado en IPFS y contratos inteligentes sobre una red BlockDAG
(traffic-storage), y un módulo de orquestación con interfaz REST
(traffic-control). Los módulos se comunican mediante APIs HTTP/JSON,
lo que garantiza bajo acoplamiento y facilita la integración con sensores reales.

La evaluación experimental demuestra que el sistema escala de forma aproximadamente lineal
con el número de sensores: para 2000 sensores, el ciclo completo de análisis y optimización
se completó en 67{,}6~s ($\approx 0{,}014$~s por sensor), con consumo de memoria estable
entre 198 y 219~MB. En cuanto al almacenamiento verificable, las transacciones alcanzaron
confirmación en la red BlockDAG en $\approx$1~s, con un costo inferior a
\$1{,}2$\times$10\textsuperscript{-5}~USD por operación, superando en más de cuatro órdenes
de magnitud el costo operativo de Ethereum y reduciendo en más del 60\% la latencia de
persistencia verificable respecto a plataformas IoT--cloud centralizadas. Los
resultados respaldan la viabilidad del sistema para despliegues municipales con recursos
computacionales limitados, incluyendo hardware embebido de bajo costo. Los hallazgos
principales fueron presentados en la 2025 IEEE CHILEAN Conference on Electrical,
Electronics Engineering, Information and Communication Technologies (CHILECON~2025).

\vspace{0.5cm}

\noindent
\textbf{Palabras Clave}: optimización de tráfico urbano, control semafórico adaptativo,
lógica difusa, optimización por enjambre de partículas, IPFS, BlockDAG, almacenamiento
verificable, microservicios, ciudad inteligente.
